%==============================================================
% RAPPORT DE PROJET - KmerPool / CitySheeP
% Transport collaboratif inter-localités au Cameroun
% Norme : ISO/IEC/IEEE
% Compilé avec : pdflatex (recommandé sur Overleaf)
%==============================================================

\documentclass[12pt, a4paper]{report}

%------ PACKAGES -----------------------------------------------
\usepackage[utf8]{inputenc}
\usepackage[T1]{fontenc}
\usepackage[french]{babel}
\usepackage{geometry}
\usepackage{graphicx}
\usepackage{float}
\usepackage{hyperref}
\usepackage{acronym}
\usepackage{glossaries}
\usepackage{array}
\usepackage{booktabs}
\usepackage{tabularx}
\usepackage{longtable}
\usepackage{xcolor}
\usepackage{fancyhdr}
\usepackage{setspace}
\usepackage{titlesec}
\usepackage{tocloft}
\usepackage{enumitem}
\usepackage{url}
\usepackage{cite}
\usepackage{listings}
\usepackage{tikz}
\usepackage{pgfplots}
\usepackage{amsmath}
\usepackage{microtype}
\usepackage{csquotes}
\usepackage{caption}
\usepackage{subcaption}

\pgfplotsset{compat=1.18}

%------ MISE EN PAGE -------------------------------------------
\geometry{
  a4paper,
  top=2.5cm, bottom=2.5cm,
  left=3cm, right=2.5cm
}
\setstretch{1.3}

%------ EN-TÊTES / PIEDS DE PAGE -------------------------------
\pagestyle{fancy}
\fancyhf{}
\fancyhead[L]{\small KmerPool -- Transport collaboratif inter-localités}
\fancyhead[R]{\small \leftmark}
\fancyfoot[C]{\thepage}
\renewcommand{\headrulewidth}{0.4pt}

%------ COULEURS -----------------------------------------------
\definecolor{primaryBlue}{RGB}{30, 80, 160}
\definecolor{accentOrange}{RGB}{220, 100, 0}
\definecolor{lightGray}{RGB}{240, 240, 240}

%------ TITRES -------------------------------------------------
\titleformat{\chapter}[block]
  {\normalfont\Large\bfseries\color{primaryBlue}}
  {\thechapter.}{12pt}{}
\titlespacing*{\chapter}{0pt}{20pt}{15pt}

\titleformat{\section}
  {\normalfont\large\bfseries\color{primaryBlue}}
  {\thesection}{10pt}{}

\titleformat{\subsection}
  {\normalfont\normalsize\bfseries}
  {\thesubsection}{8pt}{}

%------ HYPERLINKS ---------------------------------------------
\hypersetup{
  colorlinks=true,
  linkcolor=primaryBlue,
  citecolor=accentOrange,
  urlcolor=accentOrange,
  pdftitle={Rapport KmerPool},
  pdfauthor={Équipe CitySheeP},
}

%------ GLOSSAIRE ----------------------------------------------
\makeglossaries

\newglossaryentry{api}{
  name={API},
  description={Application Programming Interface -- interface de programmation permettant à deux applications de communiquer entre elles}
}
\newglossaryentry{mvc}{
  name={MVC},
  description={Model-View-Controller -- patron d'architecture logicielle séparant la logique métier, la présentation et le contrôle}
}
\newglossaryentry{rest}{
  name={REST},
  description={Representational State Transfer -- style d'architecture pour les services web}
}
\newglossaryentry{uml}{
  name={UML},
  description={Unified Modeling Language -- langage de modélisation unifié utilisé en génie logiciel}
}
\newglossaryentry{mvp}{
  name={MVP},
  description={Minimum Viable Product -- version minimale d'un produit permettant de valider une hypothèse}
}
\newglossaryentry{qr}{
  name={QR code},
  description={Quick Response Code -- code-barres 2D utilisé pour encoder des informations}
}
\newglossaryentry{sos}{
  name={SOS},
  description={Signal d'alarme international utilisé pour signaler une situation d'urgence}
}
\newglossaryentry{bd}{
  name={BD},
  description={Base de Données -- ensemble organisé d'informations structurées stockées électroniquement}
}

%------ LISTE DES ABRÉVIATIONS ---------------------------------
% (générée automatiquement, triée alphabétiquement)
\newacro{API}{Application Programming Interface}
\newacro{BD}{Base de Données}
\newacro{CSS}{Cascading Style Sheets}
\newacro{HTTP}{HyperText Transfer Protocol}
\newacro{HTTPS}{HyperText Transfer Protocol Secure}
\newacro{JSON}{JavaScript Object Notation}
\newacro{JWT}{JSON Web Token}
\newacro{MVC}{Model-View-Controller}
\newacro{MVP}{Minimum Viable Product}
\newacro{ORM}{Object-Relational Mapping}
\newacro{QR}{Quick Response}
\newacro{REST}{Representational State Transfer}
\newacro{SDK}{Software Development Kit}
\newacro{SOS}{Save Our Souls}
\newacro{SQL}{Structured Query Language}
\newacro{UI}{User Interface}
\newacro{UML}{Unified Modeling Language}
\newacro{UX}{User Experience}

%==============================================================
\begin{document}
%==============================================================

%------ PAGE DE TITRE ------------------------------------------
\begin{titlepage}
  \centering
  \vspace*{1cm}

  {\Large\bfseries Université de Yaoundé I\\[4pt]
  \normalsize Département d'Informatique}\\[2cm]

  \rule{\linewidth}{2pt}\\[0.4cm]
  {\LARGE\bfseries\color{primaryBlue}
    KmerPool -- CitySheeP\\[0.3cm]
    \large Transport collaboratif inter-localités au Cameroun}\\[0.2cm]
  \rule{\linewidth}{2pt}\\[1.5cm]

  {\large Rapport de Projet de fin d'études}\\[0.5cm]
  {\normalsize Application Mobile \& API REST}\\[2cm]

  \begin{tabular}{rl}
    \textbf{Auteurs :}    & Équipe CitySheeP \\[4pt]
    \textbf{Encadreur :}  & [Nom du professeur] \\[4pt]
    \textbf{Filière :}    & Génie Logiciel \\[4pt]
    \textbf{Année :}      & 2024 -- 2025 \\
  \end{tabular}

  \vfill
  {\small Rédigé avec \LaTeX\ -- Overleaf}
\end{titlepage}

%==============================================================
% PAGES PRÉLIMINAIRES (numérotation romaine)
%==============================================================
\pagenumbering{roman}
\setcounter{page}{1}

%------ RÉSUMÉ -------------------------------------------------
\chapter*{Résumé}
\addcontentsline{toc}{chapter}{Résumé}

KmerPool (CitySheeP) est une plateforme mobile collaborative visant à révolutionner le
transport inter-localités au Cameroun. Face à la cherté et au manque de fiabilité des
transports entre grandes villes et localités secondaires (Ntui, Obala, Bafia, etc.),
KmerPool met en relation conducteurs et passagers sur des trajets réguliers, et propose
en parallèle la livraison collaborative de colis sur ces mêmes trajets.

Le projet repose sur une architecture applicative composée d'une application mobile
multi-plateforme et d'une \ac{API} \ac{REST}. La sécurité des utilisateurs est renforcée
par un système de \ac{QR} code de vérification pré-trajet, un bouton \ac{SOS} et une
vérification d'identité. Le modèle économique s'appuie sur un abonnement mensuel,
une commission par trajet et une commission sur chaque colis livré.

Ce rapport présente l'analyse des besoins fonctionnels et non fonctionnels, la
modélisation \ac{UML}, l'architecture logique et physique, ainsi que les résultats de
l'implémentation du \ac{MVP}.

\bigskip
\noindent\textbf{Mots-clés :} transport collaboratif, covoiturage, Cameroun, application mobile,
\ac{API} \ac{REST}, livraison de colis, \ac{MVP}, sécurité, inter-localités.

%------ ABSTRACT -----------------------------------------------
\chapter*{Abstract}
\addcontentsline{toc}{chapter}{Abstract}

KmerPool (CitySheeP) is a collaborative mobile platform designed to revolutionize
inter-city transport in Cameroon. Addressing the high cost and unreliability of
transportation between major cities and secondary localities (Ntui, Obala, Bafia, etc.),
KmerPool connects drivers and passengers on regular routes while also enabling
collaborative parcel delivery along those same journeys.

The project relies on a multi-platform mobile application backed by a REST API.
User safety is enforced through a pre-trip QR code verification system, an SOS button,
and identity verification. The business model combines monthly subscriptions,
per-trip commissions, and parcel delivery commissions.

This report covers the functional and non-functional requirements analysis, UML
modelling, logical and physical architecture, and the implementation results of the MVP.

\bigskip
\noindent\textbf{Keywords:} collaborative transport, carpooling, Cameroon, mobile application,
REST API, parcel delivery, MVP, security, inter-locality.

%------ LISTE DES FIGURES (auto-générée) -----------------------
\listoffigures
\addcontentsline{toc}{chapter}{Liste des figures}

%------ LISTE DES TABLEAUX (auto-générée) ----------------------
\listoftables
\addcontentsline{toc}{chapter}{Liste des tableaux}

%------ LISTE DES ABRÉVIATIONS (auto, ordre alphabétique) ------
\chapter*{Liste des abréviations}
\addcontentsline{toc}{chapter}{Liste des abréviations}
\begin{acronym}[HTTPS]  % le plus long pour l'alignement
  \acro{API}{Application Programming Interface}
  \acro{BD}{Base de Données}
  \acro{CSS}{Cascading Style Sheets}
  \acro{HTTP}{HyperText Transfer Protocol}
  \acro{HTTPS}{HyperText Transfer Protocol Secure}
  \acro{JSON}{JavaScript Object Notation}
  \acro{JWT}{JSON Web Token}
  \acro{MVC}{Model-View-Controller}
  \acro{MVP}{Minimum Viable Product}
  \acro{ORM}{Object-Relational Mapping}
  \acro{QR}{Quick Response}
  \acro{REST}{Representational State Transfer}
  \acro{SDK}{Software Development Kit}
  \acro{SOS}{Save Our Souls}
  \acro{SQL}{Structured Query Language}
  \acro{UI}{User Interface}
  \acro{UML}{Unified Modeling Language}
  \acro{UX}{User Experience}
\end{acronym}

%------ GLOSSAIRE (auto, ordre alphabétique) -------------------
\printglossary[title=Glossaire, toctitle=Glossaire]
\addcontentsline{toc}{chapter}{Glossaire}

%------ SOMMAIRE (table des matières synthétique) --------------
\tableofcontents
\addcontentsline{toc}{chapter}{Sommaire}

%==============================================================
% CORPS DU RAPPORT (numérotation arabe)
%==============================================================
\pagenumbering{arabic}
\setcounter{page}{1}

%==============================================================
\chapter*{Introduction générale}
\addcontentsline{toc}{chapter}{Introduction générale}
%==============================================================

\section*{Contexte}
Le Cameroun dispose d'un vaste réseau routier reliant les grandes métropoles aux
localités secondaires. Pourtant, les solutions de transport structurées restent rares,
coûteuses et peu fiables pour les trajets entre villes et localités comme Ntui, Obala
ou Bafia. Des milliers de véhicules parcourent ces routes quotidiennement avec des
sièges inoccupés, tandis que des passagers et des expéditeurs de colis peinent à
trouver des options abordables et sécurisées.

\section*{Problématique}
Comment concevoir une plateforme numérique collaborative qui permette de mettre en
relation, de manière sécurisée et économique, les conducteurs et les passagers (ou
expéditeurs de colis) effectuant des trajets réguliers inter-localités au Cameroun,
dans un contexte de faible confiance entre inconnus et d'infrastructure technologique
limitée ?

\section*{Motivation}
L'absence de solutions numériques adaptées à la réalité camerounaise (connectivité
partielle, diversité linguistique, culture communautaire forte) justifie le développement
d'un outil sur mesure. KmerPool vise à valoriser les trajets existants, à réduire les
coûts de transport et à créer de la confiance via des mécanismes de notation et de
vérification d'identité.

\section*{Objectifs}

\subsection*{Objectif général}
Concevoir et développer KmerPool, une application mobile de mise en relation pour le
transport de personnes et de marchandises inter-localités au Cameroun.

\subsection*{Objectifs spécifiques}
\begin{itemize}
  \item Analyser les besoins fonctionnels et non fonctionnels des utilisateurs cibles.
  \item Modéliser le système à l'aide de diagrammes \ac{UML}.
  \item Concevoir une architecture logicielle robuste et sécurisée.
  \item Implémenter un \ac{MVP} fonctionnel (application mobile + \ac{API} \ac{REST}).
  \item Valider la solution par des tests et des captures d'écran commentées.
\end{itemize}

\section*{Plan du rapport}
Le présent rapport est organisé en trois chapitres principaux. Le \textbf{Chapitre 1}
présente les concepts clés, l'état de l'art et le bilan comparatif des solutions
existantes. Le \textbf{Chapitre 2} détaille la méthodologie, l'analyse des besoins et
la conception du système. Le \textbf{Chapitre 3} expose l'implémentation, les résultats
obtenus et le bilan technique.

%==============================================================
\chapter{Concepts, état de l'art et bilan}
%==============================================================

\section{Concepts fondamentaux}

Cette section définit les cinq concepts clés du projet KmerPool, en distinguant les
concepts métier, techniques et transversal.

\subsection{Concept métier 1 : Le covoiturage}

Le \textbf{covoiturage} désigne l'utilisation d'un véhicule privé par son conducteur pour
transporter un ou plusieurs passagers qui partagent les frais du trajet. Il se distingue
du taxi en ce que le conducteur n'est pas un professionnel du transport mais un
particulier effectuant un déplacement prévu. Dans le contexte camerounais, cette
pratique existe de manière informelle mais manque d'une infrastructure numérique
structurée pour les trajets inter-localités.

\subsection{Concept métier 2 : La livraison collaborative de colis}

La \textbf{livraison collaborative} (ou crowd-shipping) consiste à confier l'acheminement
de colis à des particuliers effectuant déjà un trajet dans la bonne direction. Ce modèle
réduit les coûts logistiques en valorisant la capacité résiduelle des véhicules en
déplacement, sans nécessiter de flotte dédiée.

\subsection{Concept technique 1 : L'application mobile multi-plateforme}

Une \textbf{application mobile multi-plateforme} est un logiciel conçu pour fonctionner
sur plusieurs systèmes d'exploitation mobiles (Android, iOS) à partir d'une base de
code unique. Des frameworks comme React Native ou Flutter permettent ce développement
croisé, réduisant les coûts et les délais par rapport au développement natif.

\subsection{Concept technique 2 : L'\ac{API} \ac{REST}}

Une \textbf{\ac{API} \ac{REST}} (Representational State Transfer) est un ensemble de
conventions architecturales permettant à des applications de communiquer via le protocole
\ac{HTTP}. Elle expose des \textit{endpoints} représentant des ressources (utilisateurs,
trajets, colis) manipulables via des méthodes standardisées (GET, POST, PUT, DELETE).
Dans KmerPool, l'\ac{API} constitue le back-end central auquel se connecte l'application
mobile.

\subsection{Concept transversal : La confiance numérique}

La \textbf{confiance numérique} est la capacité d'un système à instaurer un niveau
de sécurité et de fiabilité suffisant pour que des inconnus acceptent d'effectuer des
transactions entre eux. Elle repose sur des mécanismes tels que la vérification
d'identité, le système de notation, le \ac{QR} code de validation et le bouton \ac{SOS}.
Ce concept est transversal car il conditionne l'adoption de la plateforme par l'ensemble
des segments d'utilisateurs.

%--------------------------------------------------------------
\section{État de l'art}

Plusieurs solutions existantes abordent partiellement la problématique de KmerPool.
Cette section les recense, en distinguant les acteurs locaux (Cameroun) et
internationaux.

\subsection{Solutions au Cameroun}

À ce jour, aucune plateforme numérique structurée ne couvre le segment des trajets
inter-localités au Cameroun. Les agences de voyage (Buca Voyages, Général Express)
proposent des billets de bus mais ne permettent pas de covoiturage entre particuliers,
ni la livraison de colis via des conducteurs privés. Des groupes WhatsApp et Facebook
informels existent mais ne disposent pas d'outils de mise en relation sécurisée.

\subsection{Solutions internationales}

\begin{itemize}
  \item \textbf{BlaBlaCar} (France, Europe) : leader du covoiturage longue distance,
        ne couvre pas l'Afrique subsaharienne et ne propose pas de livraison de colis.
  \item \textbf{Uber / Bolt} : spécialisés dans le transport urbain à la demande,
        ne couvrent pas les trajets inter-localités ruraux.
  \item \textbf{Jumia Delivery} : service de livraison e-commerce limité aux grandes
        villes, sans modèle collaboratif entre particuliers.
  \item \textbf{Karos} (France) : covoiturage domicile-travail, non adapté aux longs
        trajets ruraux.
\end{itemize}

%--------------------------------------------------------------
\section{Bilan comparatif}

Le tableau~\ref{tab:bilan} présente une analyse comparative des solutions existantes
par rapport aux critères objectifs de KmerPool.

\begin{table}[H]
\centering
\caption{Tableau comparatif des solutions existantes}
\label{tab:bilan}
\renewcommand{\arraystretch}{1.3}
\begin{tabularx}{\textwidth}{|X|c|c|c|c|c|}
\hline
\textbf{Critère} & \textbf{BlaBlaCar} & \textbf{Bolt} & \textbf{Jumia Del.} & \textbf{Karos} & \textbf{KmerPool} \\
\hline
Disponible au Cameroun & Non & Partiel & Partiel & Non & \textbf{Oui} \\
\hline
Couvre les localités secondaires & Non & Non & Non & Non & \textbf{Oui} \\
\hline
Transport de personnes & Oui & Oui & Non & Oui & \textbf{Oui} \\
\hline
Livraison de colis & Non & Non & Oui & Non & \textbf{Oui} \\
\hline
Vérification d'identité & Oui & Oui & Non & Oui & \textbf{Oui} \\
\hline
Bouton SOS & Non & Oui & Non & Non & \textbf{Oui} \\
\hline
QR code de validation & Non & Non & Non & Non & \textbf{Oui} \\
\hline
Interface en français & Oui & Partiel & Oui & Oui & \textbf{Oui} \\
\hline
Gratuit (accès de base) & Oui & Non & Non & Oui & \textbf{Oui} \\
\hline
Adapté aux réseaux lents & Non & Non & Non & Non & \textbf{Oui} \\
\hline
Modèle collaboratif hybride & Non & Non & Non & Non & \textbf{Oui} \\
\hline
\end{tabularx}
\end{table}

\bigskip
\noindent\textbf{Conclusion du bilan :} Au vu de ce tableau comparatif, aucune des solutions
existantes ne remplit l'ensemble des critères objectifs identifiés. BlaBlaCar est le plus
proche fonctionnellement mais reste indisponible en Afrique subsaharienne et ne propose
pas de livraison de colis. KmerPool se positionne donc comme la seule solution répondant
à la totalité des critères dans le contexte camerounais.

%==============================================================
\chapter{Méthodologie et conception}
%==============================================================

\section{Analyse des besoins}

\subsection{Besoins fonctionnels}

Les besoins fonctionnels décrivent les actions que le système doit être capable d'effectuer.

\begin{enumerate}
  \item \textbf{Gestion des comptes utilisateurs} : inscription, connexion, vérification
        d'identité (pièce d'identité, selfie), gestion du profil.
  \item \textbf{Publication de trajets} : un conducteur peut publier un trajet avec
        les informations de départ, d'arrivée, de date, d'heure et de nombre de places.
  \item \textbf{Recherche et réservation} : un passager peut rechercher des trajets
        disponibles et effectuer une réservation.
  \item \textbf{Livraison de colis} : un expéditeur peut publier une demande de livraison
        sur un trajet existant ; un conducteur peut accepter de transporter le colis.
  \item \textbf{Messagerie intégrée} : communication directe entre conducteur, passager
        et expéditeur.
  \item \textbf{Validation QR code} : génération et scan d'un \ac{QR} code pour valider
        la prise en charge avant le départ.
  \item \textbf{Bouton SOS} : déclenchement d'une alerte d'urgence avec localisation
        GPS en temps réel.
  \item \textbf{Notation et avis} : évaluation mutuelle conducteur/passager après chaque
        trajet.
  \item \textbf{Gestion des paiements} : abonnements mensuels, commissions par trajet
        et par colis livré.
  \item \textbf{Tableau de bord} : statistiques d'utilisation pour les administrateurs.
\end{enumerate}

\subsection{Besoins non fonctionnels}

\begin{enumerate}
  \item \textbf{Performance} : temps de réponse de l'\ac{API} inférieur à 2 secondes
        pour 95\,\% des requêtes en conditions normales.
  \item \textbf{Disponibilité} : taux de disponibilité de la plateforme supérieur à
        99\,\% (SLA mensuel).
  \item \textbf{Sécurité} : chiffrement \ac{HTTPS} de toutes les communications,
        authentification par \ac{JWT}, hachage des mots de passe (bcrypt).
  \item \textbf{Scalabilité} : architecture capable de supporter une montée en charge
        progressive (de 100 à 10\,000 utilisateurs simultanés).
  \item \textbf{Accessibilité} : interface utilisable avec des connexions mobiles lentes
        (2G/3G), taille des assets optimisée.
  \item \textbf{Maintenabilité} : code documenté, couverture de tests supérieure à 70\,\%.
  \item \textbf{Internationalisation} : support du français et de l'anglais.
  \item \textbf{Conformité} : respect du Règlement Général sur la Protection des Données
        (RGPD) et de la législation camerounaise sur les données personnelles.
\end{enumerate}

%--------------------------------------------------------------
\section{Conception UML}

\subsection{Diagramme de cas d'utilisation}

La figure~\ref{fig:usecase} présente les principaux cas d'utilisation de KmerPool.
Les acteurs identifiés sont : le \textit{Passager}, le \textit{Conducteur},
l'\textit{Expéditeur} et l'\textit{Administrateur}.

\begin{figure}[H]
  \centering
  % Remplacer par le fichier image du diagramme une fois exporté depuis l'outil UML
  \fbox{\parbox{0.85\linewidth}{\centering\vspace{2cm}
    [Insérer ici le diagramme de cas d'utilisation UML -- exporter depuis draw.io ou PlantUML]
    \vspace{2cm}}}
  \caption{Diagramme de cas d'utilisation -- KmerPool}
  \label{fig:usecase}
\end{figure}

\subsection{Diagramme de classes}

Le diagramme de classes (figure~\ref{fig:classes}) modélise les entités principales
du système et leurs relations.

\begin{figure}[H]
  \centering
  \fbox{\parbox{0.85\linewidth}{\centering\vspace{2cm}
    [Insérer ici le diagramme de classes UML]
    \vspace{2cm}}}
  \caption{Diagramme de classes -- KmerPool}
  \label{fig:classes}
\end{figure}

\subsection{Diagramme d'activité : Réservation d'un trajet}

\begin{figure}[H]
  \centering
  \fbox{\parbox{0.85\linewidth}{\centering\vspace{2cm}
    [Insérer ici le diagramme d'activité pour la réservation d'un trajet]
    \vspace{2cm}}}
  \caption{Diagramme d'activité -- Réservation d'un trajet}
  \label{fig:activity}
\end{figure}

\subsection{Diagramme de séquence : Validation QR code}

\begin{figure}[H]
  \centering
  \fbox{\parbox{0.85\linewidth}{\centering\vspace{2cm}
    [Insérer ici le diagramme de séquence pour la validation QR code]
    \vspace{2cm}}}
  \caption{Diagramme de séquence -- Validation QR code}
  \label{fig:sequence}
\end{figure}

\subsection{Diagramme de déploiement}

\begin{figure}[H]
  \centering
  \fbox{\parbox{0.85\linewidth}{\centering\vspace{2cm}
    [Insérer ici le diagramme de déploiement]
    \vspace{2cm}}}
  \caption{Diagramme de déploiement -- KmerPool}
  \label{fig:deploy}
\end{figure}

%--------------------------------------------------------------
\section{Architecture du système}

\subsection{Architecture logique}

L'architecture logique de KmerPool repose sur un découpage en blocs fonctionnels
distincts, communicant via des interfaces définies. Les principaux blocs sont :

\begin{itemize}
  \item \textbf{Application mobile} (front-end) : interface utilisateur multi-plateforme
        développée avec React Native. Elle communique exclusivement avec le back-end
        via l'\ac{API} \ac{REST}.
  \item \textbf{\ac{API} \ac{REST}} (back-end) : couche de services exposant les
        ressources métier. Elle gère l'authentification \ac{JWT}, la logique métier
        et l'accès aux données.
  \item \textbf{Base de données} : stockage persistant des données (utilisateurs,
        trajets, réservations, colis, messages, notations).
  \item \textbf{Service de notifications} : envoi de notifications push en temps réel
        (réservations, alertes SOS, messages).
  \item \textbf{Service de paiement} : intégration d'un module de paiement mobile
        (Mobile Money, Orange Money).
  \item \textbf{Service cartographique} : intégration d'une \ac{API} de cartographie
        pour la géolocalisation et l'affichage des trajets.
\end{itemize}

\begin{figure}[H]
  \centering
  \fbox{\parbox{0.85\linewidth}{\centering\vspace{2cm}
    [Insérer ici le schéma d'architecture logique (blocs et flèches)]
    \vspace{2cm}}}
  \caption{Architecture logique de KmerPool}
  \label{fig:archi_logique}
\end{figure}

\subsection{Architecture physique}

L'architecture physique décrit le déploiement des composants logiciels sur
l'infrastructure matérielle et cloud.

\begin{itemize}
  \item \textbf{Serveur d'application} : instance cloud hébergeant l'\ac{API} \ac{REST}
        (ex. : Railway, Render ou AWS EC2 -- offres gratuites pour le \ac{MVP}).
  \item \textbf{Serveur de base de données} : instance PostgreSQL managée (ex. : Supabase,
        PlanetScale ou ElephantSQL).
  \item \textbf{Stockage des médias} : service de stockage objet pour les photos de profil
        et les pièces d'identité (ex. : Cloudinary, AWS S3).
  \item \textbf{CDN} : distribution des assets statiques de l'application mobile.
  \item \textbf{Appareils clients} : smartphones Android et iOS des utilisateurs finals.
\end{itemize}

\subsection{Bilan architectural}

Note : le choix du langage de programmation et des technologies spécifiques est détaillé
au chapitre suivant (Chapitre~\ref{chap:impl}), lors de la présentation de
l'implémentation.

%==============================================================
\chapter{Implémentation et résultats}
\label{chap:impl}
%==============================================================

\section{Choix technologiques}

Le tableau~\ref{tab:techno} compare les principales options technologiques évaluées
pour chaque composant de l'application, selon au moins trois aspects. Le choix final
est justifié par les courbes d'apprentissage de l'équipe et les contraintes du \ac{MVP}.

\begin{table}[H]
\centering
\caption{Comparatif des technologies -- Application mobile}
\label{tab:techno}
\renewcommand{\arraystretch}{1.3}
\begin{tabularx}{\textwidth}{|l|X|X|X|}
\hline
\textbf{Critère} & \textbf{React Native} & \textbf{Flutter} & \textbf{Natif (Java/Swift)} \\
\hline
Multi-plateforme & Oui (iOS + Android) & Oui (iOS + Android) & Non (une codebase par OS) \\
\hline
Courbe d'apprentissage & Faible (JS/TS) & Moyenne (Dart) & Élevée \\
\hline
Performance & Bonne & Très bonne & Excellente \\
\hline
Communauté & Très large & Large & Très large \\
\hline
Coût de développement & Faible & Faible & Élevé \\
\hline
\textbf{Choix} & \textbf{Retenu} & Non retenu & Non retenu \\
\hline
\end{tabularx}
\end{table}

\begin{table}[H]
\centering
\caption{Comparatif des technologies -- Back-end \ac{API}}
\label{tab:backend}
\renewcommand{\arraystretch}{1.3}
\begin{tabularx}{\textwidth}{|l|X|X|X|}
\hline
\textbf{Critère} & \textbf{Node.js / Express} & \textbf{Django REST} & \textbf{Spring Boot} \\
\hline
Langage & JavaScript / TypeScript & Python & Java \\
\hline
Performance & Bonne (event loop) & Bonne & Très bonne \\
\hline
Rapidité de dev & Très rapide & Rapide & Moyenne \\
\hline
Écosystème & Très riche & Riche & Riche \\
\hline
Déploiement gratuit & Oui (Railway, Render) & Oui & Limité \\
\hline
\textbf{Choix} & \textbf{Retenu} & Non retenu & Non retenu \\
\hline
\end{tabularx}
\end{table}

%--------------------------------------------------------------
\section{Génération de plateforme et captures d'écran}

Les sections suivantes présentent les principales interfaces de l'application KmerPool,
accompagnées de descriptions commentées.

\subsection{Écran d'accueil et authentification}

La figure~\ref{fig:screen1} présente l'écran d'accueil de KmerPool. Cet écran introduit
la proposition de valeur principale de la plateforme et offre deux points d'entrée :
la connexion pour les utilisateurs existants et l'inscription pour les nouveaux arrivants.
La charte graphique reprend les couleurs officielles du projet.

\begin{figure}[H]
  \centering
  \fbox{\parbox{0.4\linewidth}{\centering\vspace{3cm}
    [Capture d'écran 1 -- Écran d'accueil]
    \vspace{3cm}}}
  \caption{Écran d'accueil -- KmerPool}
  \label{fig:screen1}
\end{figure}

\subsection{Recherche et publication de trajets}

La figure~\ref{fig:screen2} illustre l'interface de recherche de trajets disponibles.
L'utilisateur saisit sa ville de départ, sa destination et la date souhaitée. Les résultats
affichent le conducteur, l'heure de départ, les arrêts intermédiaires et le tarif proposé.
Un système de filtres permet d'affiner la recherche par nombre de places ou par options
(colis accepté, etc.).

\begin{figure}[H]
  \centering
  \fbox{\parbox{0.4\linewidth}{\centering\vspace{3cm}
    [Capture d'écran 2 -- Recherche de trajets]
    \vspace{3cm}}}
  \caption{Interface de recherche de trajets -- KmerPool}
  \label{fig:screen2}
\end{figure}

\subsection{Validation QR code et bouton SOS}

Avant chaque trajet, le système génère un \ac{QR} code unique associé à la réservation
(figure~\ref{fig:screen3}). Le conducteur scanne le code du passager pour confirmer
sa présence, éliminant ainsi les risques de fraude. En cas d'urgence, le bouton \ac{SOS}
visible sur l'écran principal déclenche une alerte immédiate envoyant la localisation
\ac{GPS} de l'utilisateur aux contacts d'urgence définis.

\begin{figure}[H]
  \centering
  \fbox{\parbox{0.4\linewidth}{\centering\vspace{3cm}
    [Capture d'écran 3 -- Validation QR + bouton SOS]
    \vspace{3cm}}}
  \caption{Validation QR code et bouton SOS -- KmerPool}
  \label{fig:screen3}
\end{figure}

\subsection{Livraison de colis}

La figure~\ref{fig:screen4} montre l'interface de dépôt d'une demande de livraison.
L'expéditeur renseigne l'origine, la destination, le poids et la description du colis,
puis choisit parmi les conducteurs disponibles sur le trajet correspondant.

\begin{figure}[H]
  \centering
  \fbox{\parbox{0.4\linewidth}{\centering\vspace{3cm}
    [Capture d'écran 4 -- Livraison de colis]
    \vspace{3cm}}}
  \caption{Interface de livraison de colis -- KmerPool}
  \label{fig:screen4}
\end{figure}

%--------------------------------------------------------------
\section{Bilan de l'implémentation}

L'implémentation du \ac{MVP} a permis de valider les fonctionnalités core de KmerPool :
authentification sécurisée, publication et recherche de trajets, réservation, messagerie
interne, validation \ac{QR} code et bouton \ac{SOS}. Les fonctionnalités de paiement
intégré et de livraison de colis sont partiellement implémentées et feront l'objet
d'une prochaine itération.

%==============================================================
\chapter*{Conclusion générale}
\addcontentsline{toc}{chapter}{Conclusion générale}
%==============================================================

Ce projet a conduit à la conception et à l'implémentation de KmerPool, première
plateforme collaborative de transport de personnes et de marchandises inter-localités
au Cameroun. L'analyse de l'état de l'art a confirmé l'absence de solution numérique
structurée répondant aux réalités locales, justifiant ainsi la pertinence du projet.

La modélisation \ac{UML} et l'architecture \ac{REST} ont permis de poser des bases
solides pour un système scalable et maintenable. Le \ac{MVP} livré démontre la
faisabilité technique et la valeur ajoutée de la plateforme pour les trois segments
d'utilisateurs cibles : passagers, conducteurs et expéditeurs de colis.

Les perspectives d'évolution incluent l'intégration complète du paiement Mobile Money,
le déploiement d'une solution de vérification d'identité automatisée par \ac{IA},
l'extension du réseau à de nouvelles localités et le développement d'un programme de
parrainage communautaire.

%==============================================================
% ÉLÉMENTS POST-CONCLUSION (numérotation romaine)
%==============================================================
\pagenumbering{Roman}
\setcounter{page}{1}

%------ RÉFÉRENCES (standard ISO/IEC/IEEE) ---------------------
\begin{thebibliography}{99}

\bibitem{blablacar}
BlaBlaCar. \textit{How BlaBlaCar Works} [en ligne].
Disponible sur : \url{https://www.blablacar.fr} (consulté le 01/06/2025).

\bibitem{reactnative}
Facebook Inc. \textit{React Native Documentation} [en ligne].
Disponible sur : \url{https://reactnative.dev/docs/getting-started} (consulté le 01/06/2025).

\bibitem{restdissert}
Fielding, R.T., \textit{Architectural Styles and the Design of Network-based Software Architectures}.
Thèse de doctorat, University of California, Irvine, 2000.

\bibitem{umlspec}
Object Management Group. \textit{Unified Modeling Language Specification, v2.5.1}.
OMG Document formal/17-12-05, 2017.

\bibitem{isoiec}
ISO/IEC/IEEE. \textit{Systems and software engineering -- Life cycle processes --
Requirements engineering}. ISO/IEC/IEEE 29148:2018.

\bibitem{jwt}
M. Jones, J. Bradley, N. Sakimura. \textit{JSON Web Token (JWT)}.
RFC 7519, IETF, 2015. \url{https://tools.ietf.org/html/rfc7519}

\bibitem{pgdp}
Commission Nationale de l'Informatique et des Libertés (CNIL). \textit{RGPD -- Guide pratique}.
Paris : CNIL, 2022.

\bibitem{crowdship}
Devari, A., Schweizer, A.G., Mahapatra, A.
\textit{Crowdsourcing last-mile deliveries}.
Journal of Humanitarian Logistics and Supply Chain Management, vol. 7, no. 3,
pp. 259--278, 2017.

\end{thebibliography}

%------ ANNEXES -----------------------------------------------
\appendix
\chapter{Lean Canvas -- KmerPool}

\begin{table}[H]
\centering
\caption{Lean Canvas -- KmerPool v1.0}
\label{tab:leancanvas}
\renewcommand{\arraystretch}{1.5}
\begin{tabularx}{\textwidth}{|X|X|X|}
\hline
\textbf{Problème} & \textbf{Solution} & \textbf{Proposition de valeur unique} \\
Transports chers et peu fiables, véhicules avec sièges vides, livraison coûteuse entre localités, manque de confiance. &
Mise en relation conducteurs/passagers, livraison collaborative, arrêts intermédiaires, QR code, SOS, notation. &
Première plateforme hybride (personnes + colis) inter-localités au Cameroun, sécurisée et économique. \\
\hline
\textbf{Segments clients} & \textbf{Canaux} & \textbf{Sources de revenus} \\
Étudiants, travailleurs, conducteurs, commerçants, early adopters. &
WhatsApp/Facebook, parrainage, campus, marchés. &
Abonnement mensuel, commission trajet, commission colis. \\
\hline
\textbf{Structure des coûts} & \multicolumn{2}{X|}{\textbf{Avantage déloyal}} \\
Développement, hébergement cloud, maintenance, marketing. &
\multicolumn{2}{X|}{1er réseau structuré inter-localités, modèle hybride unique, communauté locale ancrée.} \\
\hline
\end{tabularx}
\end{table}

\chapter{Extraits de code source}

\lstset{
  language=JavaScript,
  basicstyle=\ttfamily\small,
  keywordstyle=\color{primaryBlue}\bfseries,
  commentstyle=\color{gray},
  stringstyle=\color{accentOrange},
  breaklines=true,
  frame=single,
  numbers=left,
  numberstyle=\tiny\color{gray},
}

\begin{lstlisting}[caption={Exemple -- Endpoint de création de trajet (Node.js / Express)}, label={lst:trajet}]
// POST /api/v1/trips
router.post('/trips', authenticate, async (req, res) => {
  try {
    const { departure, destination, date, seats, price } = req.body;
    const trip = await Trip.create({
      driverId: req.user.id,
      departure,
      destination,
      date,
      seats,
      price,
    });
    return res.status(201).json({ success: true, data: trip });
  } catch (error) {
    return res.status(500).json({ success: false, message: error.message });
  }
});
\end{lstlisting}

%------ TABLE DES MATIÈRES (fin, si référence finale exigée) ---
% (déjà générée en début de document -- commenter si doublon)

\end{document}
